\section{Future Work}
\label{sec:discussion}

\paragraph{Near term: artifact lineage and self-maintaining ecosystems.}
The most pressing near-term gap is artifact durability: like code repositories, \ara{}s decay without maintenance as dependencies rot and practices evolve, yet unmaintained artifacts are the community norm.
The natural extension is a \emph{lineage mechanism} in which each \ara{} declares its parent artifacts and expresses its contribution as a structured diff, reducing both construction cost (authors specify only the delta) and verification cost (reviewers and agents re-check only the new contribution).
Lineage also enables self-maintaining ecosystems: agents consuming an \ara{} detect and repair staleness, update deprecated dependencies, and propagate corrections upstream, so that every act of consumption becomes an act of maintenance.

\paragraph{Medium term: knowledge graph, collaborative discovery, and continuous review.}
Aggregated lineages form a queryable scientific knowledge graph that lifts collaboration and review from the document level to the corpus level.
Cross-artifact claim alignment turns literature synthesis into subgraph queries, lets reviewer agents verify that reported baselines match what cited \ara{}s recorded, and exposes trajectory conflicts where a method claimed as successful elsewhere was documented as failing.
Shared Exploration Graphs also enable collaboration formats impossible in the PDF ecosystem, from parallel continuation of open problems with documented dead ends to fine-grained attribution on live, evolving artifacts.
Review evolves in parallel: there is no single accept moment, only a claim-confidence surface that rises with replications and falls with counter-evidence, freeing human expert attention for the judgments only humans make, namely novelty, significance, and taste.

\paragraph{Long term: cross-disciplinary collective memory.}
Our evaluation is restricted to machine learning, where \ara{}'s four-layer structure aligns naturally with the dominant contribution types: algorithms, architectures, and training procedures.
Whether this structure generalizes to other disciplines remains an open question.
The Cognitive and Evidence Layers are plausibly domain-agnostic, but the Physical Layer and Exploration Graph, both premised on iterable computational experiments, may require substantial adaptation for wet-lab sciences where execution is physical rather than computational.
If these adaptations succeed, \ara{} provides a natural substrate for cross-disciplinary knowledge transfer, where documented failures in one field become actionable knowledge in another via graph traversal rather than literature search in unfamiliar notation.
